\documentclass[11pt,a4paper]{article}
% Preámbulo genérico para informes y documentos (skill informe-latex).
% Uso: copiar este archivo al proyecto como plantilla/preambulo.tex y, desde el documento,
%   \documentclass[11pt,a4paper]{article}
%   \input{plantilla/preambulo}      (ruta relativa al .tex)
\usepackage[T1]{fontenc}
\usepackage[utf8]{inputenc}
\usepackage[spanish,es-noshorthands,es-tabla]{babel}
\usepackage{lmodern}
\usepackage{microtype}
\usepackage[margin=2.3cm]{geometry}
\usepackage{xcolor}
\usepackage{booktabs}
\usepackage{tabularx}
\usepackage{array}
\usepackage{enumitem}
\usepackage{tcolorbox}
\usepackage{titlesec}
\usepackage{amsmath,amssymb}
\usepackage{graphicx}
\usepackage{caption}
\usepackage{tikz}
\usepackage[hang,flushmargin]{footmisc}
\usepackage[hidelinks]{hyperref}

\usetikzlibrary{babel,arrows.meta,positioning,calc,fit,decorations.pathreplacing}
\tcbuselibrary{skins,breakable}

\definecolor{azul}{HTML}{1F4E79}
\definecolor{gris}{HTML}{6B7280}
\definecolor{ambar}{HTML}{B45309}
\definecolor{fondoazul}{HTML}{F3F7FB}
\definecolor{fondogris}{HTML}{F4F4F5}
\definecolor{fondoambar}{HTML}{FFF7ED}

\titleformat{\section}{\Large\bfseries\color{azul}}{\thesection}{0.6em}{}
\titleformat{\subsection}{\large\bfseries}{\thesubsection}{0.6em}{}
\setcounter{tocdepth}{2}

% --- Opciones de estructura (llamar en el preámbulo del documento, antes de \begin{document}) ---
% Secciones sin número visible (siguen entrando al índice si se usan sin asterisco).
\newcommand{\sinnumeros}{\setcounter{secnumdepth}{0}}
% Cada \section empieza en página nueva.
\newcommand{\sectionbreak}{}
\newcommand{\seccionesenpaginanueva}{\renewcommand{\sectionbreak}{\clearpage}}
% Índice con hipervínculos (después de \maketitle). No fuerza página nueva entre entradas.
\newcommand{\indice}{{\renewcommand{\sectionbreak}{\addvspace{1.5em}}%
  \hypersetup{linktoc=all}\tableofcontents}}

% --- Cajas ---
% Recuadro neutro (gris): consignas, requisitos, datos de entrada, texto citado largo.
%   \begin{recuadro}{Consigna} ... \end{recuadro}
\newtcolorbox{recuadro}[1]{
  breakable, enhanced, colback=fondogris, colframe=gris,
  boxrule=0.4pt, left=6pt, right=6pt, top=4pt, bottom=4pt,
  title={\bfseries #1}, coltitle=white, colbacktitle=gris, fonttitle=\small, fontupper=\small}
% Destacado (azul): conclusión, resumen ejecutivo, definición clave.
\newtcolorbox{destacado}[1]{
  breakable, enhanced, colback=fondoazul, colframe=azul,
  boxrule=0.5pt, left=6pt, right=6pt, top=5pt, bottom=5pt,
  title={\bfseries #1}, coltitle=white, colbacktitle=azul, fonttitle=\small}
% Advertencias, supuestos, ambigüedades, erratas (ámbar).
\newtcolorbox{aviso}[1]{
  breakable, enhanced, colback=fondoambar, colframe=ambar,
  boxrule=0.5pt, left=6pt, right=6pt, top=5pt, bottom=5pt,
  title={\bfseries #1}, coltitle=white, colbacktitle=ambar, fonttitle=\small}
% Ampliación opcional (franja azul). No es una tcolorbox: las notas al pie dentro de
% una tcolorbox que se parte entre páginas se superponen con el contenido.
%   \begin{extras} ... \end{extras}      o      \begin{extras}[PARA PROFUNDIZAR] ... \end{extras}
\newenvironment{extras}[1][EXTRAS]{%
  \par\medskip\noindent
  {\setlength{\fboxsep}{3pt}\colorbox{azul}{\color{white}\bfseries\small\ #1\ }}%
  \par\nobreak\vspace{1pt}{\color{azul}\hrule height 0.6pt}\nobreak\vspace{3pt}}
 {\par\vspace{3pt}{\color{azul}\hrule height 0.6pt}\medskip}

\setlist[itemize]{itemsep=2pt, topsep=3pt, parsep=0pt}
\setlist[enumerate]{itemsep=2pt, topsep=3pt, parsep=0pt}
\captionsetup{font=small, labelfont=bf}
\setlength{\parindent}{0pt}
\setlength{\parskip}{4pt}

% --- Citas y fuentes ---
% Cita textual en el idioma original, sin traducir.
\newcommand{\cita}[1]{``\textit{#1}''}
% Fuente en nota al pie (nunca en el cuerpo):  \fuente{Autor, \textit{Título}, §2.1, p.~34}
% Varias fuentes en una nota: separarlas con «;».
\newcommand{\fuente}[1]{\footnote{#1.}}
% Nota inicial que explica abreviaturas de fuentes:  \notafuentes{\textbf{AIMA}: Russell \& Norvig, ...}
\newcommand{\notafuentes}[1]{{\small\color{gris} Las fuentes van en notas al pie. #1\par}}

% Íconos reutilizables (pics de TikZ). Uso: \pic[scale=0.8] at (x,y) {robot};
\tikzset{
  pics/robot/.style={code={
    \draw[thick, fill=azul!15] (-0.35,0.15) rectangle (0.35,0.65);            % cabeza
    \fill[azul] (-0.15,0.42) circle (0.06) (0.15,0.42) circle (0.06);         % ojos
    \draw[thick] (0,0.65) -- (0,0.8); \fill[ambar] (0,0.83) circle (0.05);    % antena
    \draw[thick, fill=azul!15] (-0.45,-0.45) rectangle (0.45,0.12);           % cuerpo
    \draw[thick] (-0.45,-0.1) -- (-0.65,0.1) (0.45,-0.1) -- (0.65,0.1);       % brazos
    \draw[thick, fill=gris!40] (-0.25,-0.55) circle (0.12) (0.25,-0.55) circle (0.12); % ruedas
  }},
  pics/compu/.style={code={
    \draw[thick, fill=azul!15, rounded corners=1pt] (-0.55,-0.25) rectangle (0.55,0.5); % pantalla
    \node[font=\tiny\ttfamily\bfseries, text=azul] at (0,0.12) {</>};
    \draw[thick] (0,-0.25) -- (0,-0.4); \draw[thick] (-0.25,-0.4) -- (0.25,-0.4);       % pie
  }},
  pics/persona/.style={code={
    \draw[thick] (0,0.35) circle (0.1);
    \draw[thick] (0,0.25) -- (0,-0.05) (-0.15,0.15) -- (0.15,0.15)
                 (0,-0.05) -- (-0.12,-0.3) (0,-0.05) -- (0.12,-0.3);
  }},
  pics/arbol/.style={code={
    \draw[thick, brown!60!black] (0,-0.3) -- (0,0);
    \fill[green!50!black!70] (0,0.15) circle (0.2);
  }},
  pics/documento/.style={code={
    \draw[thick, fill=white] (-0.2,-0.28) -- (-0.2,0.28) -- (0.1,0.28) -- (0.2,0.18) -- (0.2,-0.28) -- cycle;
    \foreach \y in {0.1,0,-0.1} \draw[gris] (-0.12,\y) -- (0.12,\y);
  }},
  pics/sobre/.style={code={
    \draw[thick, fill=white] (-0.25,-0.17) rectangle (0.25,0.17);
    \draw[thick] (-0.25,0.17) -- (0,-0.02) -- (0.25,0.17);
  }},
}

% Opcionales:  \sinnumeros   \seccionesenpaginanueva

\title{\vspace{-1.5cm}\textbf{Título del informe}\\[3pt]
\large Subtítulo}
\author{\normalsize Autor --- Institución --- 2026}
\date{}

\begin{document}
\maketitle
\vspace{-1cm}
\notafuentes{\textbf{AIMA}: Russell \& Norvig, \textit{Artificial Intelligence: A Modern Approach}, 3.ª ed. (2010).}

\indice

\section{Introducción}

\begin{recuadro}{Objetivo}
Qué se pide o qué se busca responder.
\end{recuadro}

\begin{itemize}
  \item Primer punto, con su fuente en nota al pie.\fuente{AIMA, §2.1, p.~34}
  \item Una cita en el idioma original: \cita{An agent is anything that can be viewed as perceiving its environment}.
  \item Un porcentaje en matemática: $35\%$; un decimal: $0{,}65$.
\end{itemize}

\begin{center}
\small
\begin{tabularx}{\linewidth}{@{}>{\bfseries}l X X@{}}
\toprule
 & \textbf{Opción A} & \textbf{Opción B} \\
\midrule
Costo & bajo & alto \\
Ventaja & simple & flexible \\
\bottomrule
\end{tabularx}
\end{center}

\begin{figure}[h]
\centering
\begin{tikzpicture}[flecha/.style={-{Stealth[length=2.2mm]}, thick, draw=azul}]
  \pic at (0,0) {persona};
  \pic at (3,0) {compu};
  \draw[flecha] (0.4,0) -- node[above, font=\scriptsize] {pide} (2.3,0);
\end{tikzpicture}
\caption{Pictograma con íconos del preámbulo.}
\end{figure}

\begin{extras}
\begin{itemize}
  \item Información nueva que amplía, pero no hace falta para entender lo de arriba.
\end{itemize}
\end{extras}

\section{Conclusión}

\begin{destacado}{Conclusión}
La idea principal en una o dos líneas.
\end{destacado}

\begin{aviso}{Supuesto}
Algo que hubo que suponer porque faltaba un dato.
\end{aviso}

\end{document}
